% Part: computability
% Chapter: computability-theory
% Section: equiv-ce-defs

\documentclass[../../../include/open-logic-section]{subfiles}

\begin{document}

\olfileid{cmp}{thy}{eqc}

\olsection[قابلِ حساب فہرست بنن جوگ سیٹاں دیاں تعریفاں]{قابلِ حساب
  گنتی وار فہرست بنن جوگ سیٹاں دیاں ہم ارز تعریفاں}


ہِٹھاں دتے بیان ایہ سمجھاؤن دے کئی اہم تے ہم ارز
طریقے نیں کہ سیٹ دے قابلِ حساب گنتی وار فہرست
بنن جوگ ہون دا مطلب کی اے۔

\begin{thm}
\ollabel{thm:ce-equiv}
من لو $S$ قدرتی عدداں دا سیٹ اے۔ تاں ہِٹھلیاں
شرطاں آپس وچ ہم ارز نیں:
\begin{enumerate}
\item\ollabel{case:ce} $S$ قابلِ حساب گنتی وار فہرست بنن جوگ اے۔
\item\ollabel{case:ran-pc} $S$ کسے \emph{جزوی} قابلِ حساب فنکشن دیاں حاصل قدراں دا سیٹ اے۔
\item\ollabel{case:ran-prim} $S$ خالی اے یا کسے بنیادی بازگشتی فنکشن دیاں حاصل قدراں دا سیٹ اے۔
\item\ollabel{case:ce-domain} $S$ کسے جزوی قابلِ حساب فنکشن دا \emph{دائرۂ تعریف} اے۔
\end{enumerate}
\end{thm}

\begin{explain}
پہلیاں تِن شرطاں دا مطلب اے کہ کسے وی غیر خالی
قابلِ حساب گنتی وار فہرست بنن جوگ سیٹ دی گنتی وار
فہرست قابلِ حساب فنکشن، جزوی قابلِ حساب فنکشن یا
بنیادی بازگشتی فنکشن نال بنائی جا سکدی اے۔ چوتھی
شرط دسدی اے کہ جے $S$ قابلِ حساب گنتی وار فہرست
بنن جوگ اے، تاں کسے اشاریے~$e$ لئی
\[
S = \Setabs{x}{\cfind{e}(x) \fdefined}.
\]
ہور لفظاں وچ، $S$ اوہناں انپٹاں دا سیٹ اے جنہاں
لئی $\cfind{e}$ دی حسابی کارروائی رُک جاندی اے۔
ایس واسطے قابلِ حساب گنتی وار فہرست بنن جوگ سیٹاں
نوں کدی کدی \emph{نیم فیصلہ کرن جوگ
(semi-decidable)} وی آکھیا جاندا اے: جے کوئی عدد
سیٹ وچ ہووے تاں آخرکار ``ہاں'' مل جاندی اے، پر
جے اوہ سیٹ وچ نہ ہووے تاں ``نہیں'' کدی نہیں ملدی!{}
\end{explain}

\begin{proof}
ہر بنیادی بازگشتی فنکشن قابلِ حساب ہوندا اے تے ہر
قابلِ حساب فنکشن جزوی قابلِ حساب وی ہوندا اے، ایس
لئی \olref{case:ran-prim} توں \olref{case:ce}
نکلدا اے، تے \olref{case:ce} توں
\olref{case:ran-pc} نکلدا اے۔ (دھیان دیو: جے $S$
خالی ہووے، تاں $S$ اوس جزوی قابلِ حساب فنکشن
دیاں حاصل قدراں دا سیٹ اے جو کسے وی انپٹ اُتے
تعریف شدہ نہیں۔) جے اسی ایہ وکھا دئیے کہ
\olref{case:ran-pc} توں \olref{case:ran-prim}
نکلدا اے، تاں پہلیاں تِن شرطاں دی ہم ارزیت ثابت
ہو جائے گی۔

ہن من لو $S$ جزوی قابلِ حساب فنکشن
$\cfind{e}$ دیاں حاصل قدراں دا سیٹ اے۔ جے $S$
خالی اے تاں گل ثابت ہو گئی۔ نہیں تاں $a$ نوں
$S$ دا کوئی وی رکن چن لو۔ کلینی دے نارمل روپ
والے قضیے مطابق اسی لکھ سکدے آں:
\[
\cfind{e}(x) = U(\umin{s}{T(e, x, s)}).
\]
خاص طور اُتے، $\cfind{e}(x) \fdefined$ تے $= y$
اُدوں تے صرف اُدوں ہوندا اے جدوں کوئی $s$ ہووے
جس لئی $T(e, x, s)$ درست ہووے تے $U(s) = y$
ہووے۔ $f(z)$ دی تعریف ایویں کرو:
\[
f(z) = \begin{cases}
  U((z)_1) & \text{جے $T(e, (z)_0, (z)_1)$} \\
  a        & \text{ہور صورت وچ۔}
\end{cases}
\]
تاں $f$ بنیادی بازگشتی اے، کیوں جے $T$ تے $U$
بنیادی بازگشتی نیں۔ \iftag{TMs}{ٹیورنگ مشیناں
دے لحاظ نال، جے $z$ جوڑی
$\tuple{(z)_0, (z)_1}$ دا کوڈ ہووے تے $(z)_1$،
مشین~$M_e$ دی انپٹ $(z)_0$ اُتے رُک جان والی
حسابی کارروائی ہووے، تاں $f$ اوس کارروائی دا
آؤٹ پٹ دیندا اے؛ نہیں تاں اوہ~$a$ واپس دیندا اے۔}

ساڈے لئی ایہ ثابت کرنا باقی اے کہ $S$ عین~$f$
دیاں حاصل قدراں دا سیٹ اے؛ یعنی ہر قدرتی
عدد~$y$ لئی، $y \in S$ اُدوں تے صرف اُدوں
جدوں اوہ~$f$ دیاں حاصل قدراں وچ ہووے۔ پہلے
رخ وچ من لو $y \in S$۔ تاں $y$،
$\cfind{e}$ دیاں حاصل قدراں وچ اے؛ ایس لئی
کسے $x$ تے~$s$ لئی $T(e,x,s)$ درست اے
تے $U(s) = y$۔ پر فیر $y = f(\tuple{x,s})$۔
اُلٹے رخ وچ من لو $y$، فنکشن~$f$ دیاں
حاصل قدراں وچ اے۔ تاں یا $y = a$ اے،
یا کسے~$z$ لئی $T(e,(z)_0,(z)_1)$ تے
$U((z)_1) = y$۔ % PNB-CMP-015: latter case uses input (z)_0, not unbound x.
دوجی صورت وچ $\cfind{e}((z)_0) \fdefined =
y$ ہون کرکے، دوناں صورتاں وچ $y$،
$S$ وچ اے۔

(علامت $\cfind{e}(x) \fdefined = y$ دا مطلب
اے کہ ``$\cfind{e}(x)$ تعریف شدہ اے تے
$y$ دے برابر اے۔'' اسی $\cfind{e}(x) =
y$ وی لکھ سکدے ساں، پر وادھو تیر کدی کدی
ایہ یاد کران لئی فائدہ مند اے کہ ایتھے فنکشن
جزوی اے۔)

\olref{thm:ce-equiv} دا ثبوت پورا کرن لئی
ایہ وکھانا کافی اے کہ \olref{case:ce} تے~
\olref{case:ce-domain} ہم ارز نیں۔ پہلاں
وکھائیے کہ \olref{case:ce} توں~
\olref{case:ce-domain} نکلدا اے۔ % PNB-CMP-016: source omits empty CE case in this direction.
جے زیرِ غور سیٹ خالی اے تاں اوہ کسے وی انپٹ
اُتے تعریف نہ ہون والے جزوی قابلِ حساب فنکشن
دا دائرۂ تعریف اے۔ ہن من لو $S$
قابلِ حساب فنکشن~$f$ دیاں حاصل قدراں دا سیٹ
اے، یعنی
\[
S = \Setabs{y}{\text{کسے $x$ لئی، } f(x) = y}.
\]
ایہ تعریف کرو:
\[
g(y) = \umin{x}{(f(x) = y)}.
\]
تاں $g$ جزوی قابلِ حساب فنکشن اے، تے $g(y)$
اُدوں تے صرف اُدوں تعریف شدہ اے جدوں
کسے~$x$ لئی $f(x) = y$ ہووے۔ ہور
لفظاں وچ، $g$ دا دائرۂ تعریف~$f$
دیاں حاصل قدراں دا سیٹ اے۔
\iftag{TMs}{ٹیورنگ مشیناں دے لحاظ نال: جے
ٹیورنگ مشین~$F$، $S$ دے رکناں دی
گنتی وار فہرست بناندی اے، تاں مشین $G$
نوں ایداں بناؤ کہ اوہ $S$ بارے نیم فیصلہ کرے۔
اوہ~$F$ دے آؤٹ پٹاں وچ دِتّا رکن لبھدی
رہندی اے؛ رکن لبھ جائے تاں رُک جاندی اے،
نہیں تاں ہمیشہ تلاش جاری رکھدی اے۔}{}

آخر وچ \olref{case:ce-domain} توں~
\olref{case:ce} نکلنا وکھاؤن لئی من لو
$S$ جزوی قابلِ حساب فنکشن~$\cfind{e}$
دا دائرۂ تعریف اے، یعنی
\[
S = \Setabs{x}{\cfind{e}(x) \fdefined}.
\]
جے $S$ خالی اے تاں گل ثابت اے؛ نہیں تاں
$a$ نوں~$S$ دا کوئی وی رکن چن لو۔
$f$ دی تعریف ایویں کرو:
\[
f(z) = \begin{cases}
(z)_0 & \text{جے $T(e,(z)_0,(z)_1)$} \\
a & \text{ہور صورت وچ۔}
\end{cases}
\]
تاں اُتّے وانگ کوئی عدد $x$، فنکشن~$f$
دیاں حاصل قدراں وچ اُدوں تے صرف اُدوں
اے جدوں $\cfind{e}(x) \fdefined$، یعنی
اُدوں تے صرف اُدوں جدوں $x \in S$۔
\iftag{TMs}{ٹیورنگ مشیناں دے لحاظ نال:
جے مشین $M_e$، $S$ بارے نیم فیصلہ
کردی اے، تاں ٹیورنگ مشیناں دیاں ساریاں ممکنہ
حسابی کارروائیاں دی فہرست نوں اک اک کرکے ویکھو تے
اوہناں رُک جان والیاں کارروائیاں دے
متعلقہ انپٹ واپس دیو؛ ایس طرح $S$
دے رکناں دی فہرست بن جاندی اے۔}{}
\end{proof}

\olref{thm:ce-equiv} دی شرط~
\olref{case:ce-domain}، قابلِ حساب گنتی وار
فہرست بنن جوگ سیٹاں دی فہرست بناؤن دا اک
سَوکھا طریقہ دیندی اے: ہر~$e$ لئی $W_e$
نوں $\cfind{e}$ دا دائرۂ تعریف من لو،
یعنی
\[
W_e = \Setabs{x}{\cfind{e}(x) \fdefined}.
\]
تاں جے $A$ کوئی وی قابلِ حساب گنتی وار
فہرست بنن جوگ سیٹ اے، تاں $A = W_e$
کسے~$e$ لئی ہوندا اے۔

اگلا قضیہ قابلِ حساب گنتی وار فہرست بنن
جوگ سیٹاں دی اک ہور تعریف دیندا اے۔

\begin{thm}
\ollabel{thm:exists-char}
سیٹ $S$ اُدوں تے صرف اُدوں قابلِ حساب
گنتی وار فہرست بنن جوگ اے جدوں کوئی
قابلِ حساب تعلق $R(x,y)$ ہووے جس لئی
\[
S = \Setabs{ x }{ \lexists[y][R(x,y)] }.
\]
\end{thm}

\begin{proof}
پہلے رخ وچ من لو $S$ قابلِ حساب گنتی وار
فہرست بنن جوگ اے۔ تاں کسے $e$ لئی،
$S = W_e$۔ ایس~$e$ لئی اسی $S$
نوں ایویں لکھ سکدے آں:
\[
S = \Setabs{ x }{ \lexists[y][T(e, x, y)] }.
\]
اُلٹے رخ وچ من لو $S = \Setabs{ x }{
  \lexists[y][R(x, y)] }$۔ $f$~دی
تعریف ایویں کرو:
\[
f(x) \simeq \umin{y}{R(x, y)}.
\]
تاں $f$ جزوی قابلِ حساب اے، تے $S$
فنکشن~$f$ دا دائرۂ تعریف اے۔
\end{proof}


\end{document}
